\chapter{Planteamiento del problema}

La interpretaci\'on de resonancias magn\'eticas lumbares requiere revisar secuencias, identificar estructuras anat\'omicas, comparar niveles y, cuando corresponde, realizar mediciones. Aunque los visores m\'edicos actuales permiten observar im\'agenes y efectuar mediciones manuales, la delimitaci\'on de estructuras y la organizaci\'on de resultados cuantitativos pueden seguir dependiendo de tareas manuales o semimanuales. En RM lumbar sagital, estructuras como cuerpos vertebrales, discos intervertebrales y canal espinal son relevantes para el an\'alisis estructural. Su delimitaci\'on manual puede consumir tiempo, variar entre observadores y generar resultados dif\'iciles de reutilizar si no quedan asociados a una estructura, un nivel anat\'omico y una imagen de origen. Esta situaci\'on afecta especialmente la trazabilidad de mediciones y observaciones cuando se trabaja con capturas, registros dispersos o informes redactados sin una salida estructurada.

El problema central que aborda este proyecto puede formularse de la siguiente manera:

\begin{quote}
\itshape
Existe una oportunidad de asistencia inform\'atica en el an\'alisis estructural de resonancias magn\'eticas lumbares sagitales, dado que la segmentaci\'on de estructuras anat\'omicas, la obtenci\'on de mediciones simples y la generaci\'on de salidas revisables a\'un pueden requerir tareas manuales, variables y poco integradas dentro de un mismo flujo funcional.
\end{quote}

La dificultad no se ubica en reemplazar la interpretaci\'on profesional, sino en asistir tareas operativas espec\'ificas. Un sistema de apoyo puede proponer segmentaciones, calcular mediciones geom\'etricas y presentar resultados visuales, pero debe permitir revisi\'on, correcci\'on o descarte por parte del usuario. Desde una perspectiva de ingenier\'ia, el desaf\'io consiste en integrar carga de im\'agenes, preprocesamiento, segmentaci\'on autom\'atica, mediciones, visualizaci\'on superpuesta y salida editable en un flujo coherente. La soluci\'on debe ser trazable, reproducible y acotada al contexto acad\'emico, sin presentarse como producto cl\'inico certificado ni como sistema de diagn\'ostico aut\'onomo.
